\section{How to use this file}
\label{sec:orgba55d0f}
This file shows how Org marks up text: emphasis, sub- and superscripts,
special symbols, blocks, comments, line breaks, keywords and macros. Most of
this markup is \textbf{shown} by the editor (colours, hidden markers) and \textbf{used}
by the exporter (HTML, \LaTeX{}, plain text\ldots{}).

Start Neovim from the root of the repository with the bundled config:

\begin{verbatim}
nvim -u examples/minimal_init.lua examples/02-markup.org
\end{verbatim}

\begin{itemize}
\item The file opens folded. \texttt{<Tab>} on a heading opens it, \texttt{<S-Tab>} cycles
the whole buffer.
\item \texttt{<prefix>} means \texttt{<leader>o} (\texttt{<Space>o} with the bundled config). \texttt{g?}
lists every key of the buffer.
\item \texttt{u} undoes; \texttt{git checkout examples/02-markup.org} restores the file.
\item \textbf{Try:} lines are exercises, \textbf{Expect:} lines say what you should see.
\item Lines starting with =\# = are comments: notes for you, never exported.
\end{itemize}
\subsection{How to check the result: export one subtree}
\label{sec:org3ef3b96}
Markup is easiest to verify by exporting. With the cursor anywhere inside a
heading, this command exports \textbf{only that subtree}, without the HTML page
around it, into a scratch buffer:

\begin{verbatim}
:Org export html subtree body buffer
\end{verbatim}

A split named "Org HTML Export" opens with the HTML. Close it with \texttt{:q}.
For a plain-text view use \texttt{utf8} instead of \texttt{html}. The same exports are in
the menu of \texttt{<prefix>e} (Emacs \texttt{<C-c><C-e>}): press \texttt{s} (subtree only) and
\texttt{b} (body only), then \texttt{h H} (HTML buffer) or \texttt{t U} (UTF-8 buffer). Export
itself is covered in \href{19-export.org}{19-export.org}.

Most \textbf{Expect:} lines below quote the exact HTML you should find.
\subsection{Keys in this file}
\label{sec:orga0d449c}
\begin{center}
\begin{tabular}{ll}
Key & What it does\\
\hline
\texttt{<prefix>E} & emphasize: wrap the selection in \texttt{*} \texttt{/} \texttt{\_}\ldots{}\\
\texttt{<C-c><C-x><C-f>} & the same, Emacs key\\
\texttt{<C-c><C-x>\textbackslash{}} & toggle pretty entities (\texttt{\textbackslash{}alpha} shown as α)\\
\texttt{<C-c>:} & toggle fixed-width =: = on the line / selection\\
\texttt{<prefix>ib} & insert a block (quote, example, src\ldots{})\\
\texttt{<prefix>hC} & toggle the COMMENT keyword of a headline\\
\texttt{<Tab>} & on a \texttt{\#+begin\_} line: fold / unfold the block\\
\texttt{<prefix>e} & export menu\\
\texttt{<prefix>xl} & preview \LaTeX{} fragments as images\\
\end{tabular}
\end{center}
\section{Emphasis}
\label{sec:org9627784}
Six markers change how text looks. Put the marker right before the first
character and right after the last one:

\begin{center}
\begin{tabular}{lll}
You write & You get & HTML\\
\hline
\texttt{*bold*} & \textbf{bold} & \texttt{<b>bold</b>}\\
\texttt{/italic/} & \emph{italic} & \texttt{<i>italic</i>}\\
\texttt{\_underlined\_} & \uline{underlined} & \texttt{<span class}"underline">\ldots{}=\\
\texttt{+strike+} & \sout{strike} & \texttt{<del>strike</del>}\\
\texttt{=verbatim=} & \texttt{verbatim} & \texttt{<code>verbatim</code>}\\
\texttt{\textasciitilde{}code\textasciitilde{}} & \texttt{code} & \texttt{<code>code</code>}\\
\end{tabular}
\end{center}

\texttt{verbatim} and \texttt{code} look the same in most exports; the difference is
that nothing inside them is interpreted: no emphasis, no links, no
entities. Use \texttt{code} for code and \texttt{verbatim} for keys, file names and
literal strings (this is the convention of the Org manual).

Emphasis can be nested, and can span \textbf{two} lines, but not more.

\textbf{Try:} open "Emphasis examples" below and look at the colours: each kind of
markup has its own highlight group (\texttt{OrgBold}, \texttt{OrgItalic}, \texttt{OrgUnderline},
\texttt{OrgStrikethrough}, \texttt{OrgVerbatim}, \texttt{OrgCode}). Then, with the cursor in it,
run \texttt{:Org export html subtree body buffer}.

\textbf{Expect:} the HTML contains, one per line:
\begin{verbatim}
<b>bold</b>, <i>italic</i>, <span class="underline">underlined</span>,
<del>strike-through</del>, <code>verbatim</code> and <code>code</code>.
<b>bold with <i>italic inside</i> it</b>
<code>*not bold*</code>
<b>emphasis over
two lines</b>
\end{verbatim}
\subsection{Emphasis examples}
\label{sec:org2dc5823}
\textbf{bold}, \emph{italic}, \uline{underlined},
\sout{strike-through}, \texttt{verbatim} and \texttt{code}.
\textbf{bold with \emph{italic inside} it}
\texttt{*not bold*} because verbatim shows its text as it is.
\textbf{emphasis over
two lines} works.
\subsection{When emphasis does not work}
\label{sec:orgeab1bb9}
The markers are only recognised in the right context (Org's
\texttt{org-emphasis-regexp-components}):

\begin{itemize}
\item The opening marker must be at the start of a line or after a space or one
of \texttt{-('"\{}. The closing marker must be followed by a space, the end of the
line, or one of \texttt{-.,:!?;'")\}\textbackslash{}[}. So markers inside a word do nothing:
\texttt{a*b*c} stays as it is.
\item The text inside cannot start or end with a space: \texttt{* no *} is not bold.
\item Emphasis spans at most two lines. Three lines: no emphasis.
\item Inside \texttt{verbatim} and \texttt{code} nothing is interpreted.
\end{itemize}

\textbf{Try:} open "Things that are not emphasis" and export it as above.

\textbf{Expect:} none of these lines has \texttt{<b>}, \texttt{<i>} or \texttt{<del>} in the HTML:
the stars, slashes and plus signs come out as plain characters:
\begin{verbatim}
a*b*c in the middle of a word.
Here * no stars * because of the spaces inside.
2+3+4 is just arithmetic.
path/to/file is not italic.
\end{verbatim}
(Underscores are different: \texttt{snake\_case} gives a subscript, see the next
topic.)
\subsection{Things that are not emphasis}
\label{sec:orgaf6c356}
a*b*c in the middle of a word.
Here * no stars * because of the spaces inside.
2+3+4 is just arithmetic.
path/to/file is not italic.
\subsection{Adding emphasis with a key}
\label{sec:org3030266}
\texttt{<prefix>E} (Emacs \texttt{<C-c><C-x><C-f>}) asks for a marker (\texttt{*} \texttt{/} \texttt{\_} \texttt{+}
\texttt{=} \texttt{\textasciitilde{}}) and:
\begin{itemize}
\item in Visual mode, wraps the selection in it; if the selection already has a
marker, it is replaced; answering \texttt{<Space>} removes the markers;
\item in Normal mode, inserts a pair of markers and puts you in Insert mode
between them.
\end{itemize}
Spaces are added when needed so the markers are recognised.

\textbf{Try:} on the line "Make this word bold." below, put the cursor on "word",
press \texttt{viw}, then \texttt{<prefix>E} and \texttt{*}.

\textbf{Expect:} the line reads \texttt{Make this *word* bold.} and "word" turns bold.

\textbf{Try:} put the cursor on the first star of \texttt{*word*}, select up to the
second star with \texttt{vf*}, press \texttt{<prefix>E} and \texttt{/}.

\textbf{Expect:} the stars are replaced: \texttt{Make this /word/ bold.}

\textbf{Try:} on the line "Type between alpha beta", put the cursor on the \texttt{b} of
"beta", press \texttt{<prefix>E} and \texttt{\textasciitilde{}}, type \texttt{ls} and press \texttt{<Esc>}.

\textbf{Expect:} \texttt{Type between alpha \textasciitilde{}ls\textasciitilde{} beta}: the pair was inserted before
"beta" and a space was added after it, so the markers are recognised.
\subsubsection{Emphasize playground}
\label{sec:org28d6e92}
Make this word bold.
Type between alpha beta
\subsection{Hiding the markers}
\label{sec:org8463113}
With \texttt{ui.hide\_emphasis\_markers = true} the markers are concealed: you see
\textbf{bold} without its stars. The markers stay in the file, and show again on
the line under the cursor (depending on \texttt{'concealcursor'}).

\textbf{Try:} run these two commands, then look at "Emphasis examples":

\begin{verbatim}
:lua require("org.config").opts.ui.hide_emphasis_markers = true
:w | e
\end{verbatim}

\textbf{Expect:} the markers disappear from lines away from the cursor. Set it
back to \texttt{false} the same way.
\section{Subscripts and superscripts}
\label{sec:org1183b91}
\texttt{\_} makes a subscript and \texttt{\textasciicircum{}} a superscript: \texttt{x\textasciicircum{}2}, \texttt{a\_i}. Without braces
the script extends over all the letters and digits that follow, so write
\texttt{H\_\{2\}O} (not \texttt{H\_2O}, which subscripts "2O"). Braces are always safe:
\texttt{x\textasciicircum{}\{10\}}, \texttt{a\_\{ij\}}.

Two settings control them:
\begin{itemize}
\item In the buffer, with \texttt{ui.pretty\_entities} on (see Entities below), scripts
are drawn with Unicode super/subscript characters when all of their
characters have one (x², aᵢⱼ), else highlighted with \texttt{OrgSuperscript} /
\texttt{OrgSubscript}.
\item In the export, \texttt{\#+OPTIONS: \textasciicircum{}:\{\}} (or \texttt{:EXPORT\_OPTIONS: \textasciicircum{}:\{\}} on a
subtree) only accepts scripts in braces, so \texttt{snake\_case} stays intact.
\texttt{\textasciicircum{}:nil} turns them off.
\end{itemize}

\textbf{Try:} export "Scripts with defaults" (\texttt{:Org export html subtree body
buffer}).

\textbf{Expect:}
\begin{verbatim}
H<sub>2</sub>O and E=mc<sup>2</sup> and a<sub>ij</sub>
H<sub>2O</sub> is wrong: "2O" is subscripted.
snake<sub>case</sub>
\end{verbatim}

\textbf{Try:} export "Scripts with braces only" the same way.

\textbf{Expect:} \texttt{H\_2O} and \texttt{snake\_case} stay as they are, only the braced ones
change:
\begin{verbatim}
H_2O snake_case x<sup>2</sup> a<sub>ij</sub>
\end{verbatim}
\subsection{Scripts with defaults}
\label{sec:orgbdb4163}
H\textsubscript{2}O and E=mc\textsuperscript{2} and a\textsubscript{ij}
H\textsubscript{2O} is wrong: "2O" is subscripted.
snake\textsubscript{case}
\subsection{Scripts with braces only}
\label{sec:org6852f3b}
H\textsubscript{2O} snake\textsubscript{case} x\textsuperscript{2} a\textsubscript{ij}
\section{Entities}
\label{sec:org3d780a4}
Entities are \TeX{}-like names for symbols: \texttt{\textbackslash{}alpha}, \texttt{\textbackslash{}to}, \texttt{\textbackslash{}deg},
\texttt{\textbackslash{}copy}\ldots{} Org knows hundreds of them (the same list as Emacs'
\texttt{org-entities}). Each exporter writes the right thing: \texttt{\&alpha;} in HTML,
\texttt{\textbackslash{}alpha} in \LaTeX{}, \texttt{α} in UTF-8 text. End an entity with \texttt{\{\}} when a letter
follows: \texttt{\textbackslash{}alpha\{\}beta}.

In the buffer, \texttt{ui.pretty\_entities = true} (or \texttt{\#+STARTUP: entitiespretty})
shows them as their Unicode character. \texttt{<C-c><C-x>\textbackslash{}} toggles this for the
current buffer. Entities are never prettified in blocks, code, verbatim or
links. As in Emacs, only entities with a one-character symbol are drawn:
\texttt{\textbackslash{}sin} or \texttt{\textbackslash{}lim} stay as written in the buffer (they still export as
\texttt{sin} and \texttt{lim}).

\begin{center}
\begin{tabular}{llll}
Write & Shows as & Write & Shows as\\
\hline
\texttt{\textbackslash{}alpha} & α & \texttt{\textbackslash{}to} & →\\
\texttt{\textbackslash{}beta} & β & \texttt{\textbackslash{}larr} & ←\\
\texttt{\textbackslash{}pi} & π & \texttt{\textbackslash{}deg} & °\\
\texttt{\textbackslash{}infty} & ∞ & \texttt{\textbackslash{}copy} & ©\\
\texttt{\textbackslash{}pm} & ± & \texttt{\textbackslash{}euro} & €\\
\texttt{\textbackslash{}times} & × & \texttt{\textbackslash{}nbsp} & (space)\\
\end{tabular}
\end{center}

\textbf{Try:} open "Entity examples" and press \texttt{<C-c><C-x>\textbackslash{}}.

\textbf{Expect:} \texttt{\textbackslash{}alpha} is drawn as α, \texttt{\textbackslash{}to} as →, \texttt{\textbackslash{}deg\{\}} as °, and the
scripts \texttt{x\textasciicircum{}2} and \texttt{a\_\{ij\}} as x² and aᵢⱼ. Press \texttt{<C-c><C-x>\textbackslash{}} again to see
the plain text.

\textbf{Try:} export "Entity examples" to HTML.

\textbf{Expect:}
\begin{verbatim}
&alpha; &beta; &pi; &rarr; 20&deg;C &copy; 2026
&alpha;beta (the {} ends the name)
\end{verbatim}
\subsection{Entity examples}
\label{sec:orgc4eb537}
\(\alpha\) \(\beta\) \(\pi\) \(\to\) 20\textdegree{}C \textcopyright{} 2026
\(\alpha\)beta (the \{\} ends the name)
x\textsuperscript{2} and a\textsubscript{ij}
\section{Special strings and smart quotes}
\label{sec:org5570f68}
When exporting, a few character sequences become typographic symbols
(\texttt{\#+OPTIONS: -:t}, the default; \texttt{-:nil} turns it off):

\begin{center}
\begin{tabular}{lll}
Write & Becomes & HTML\\
\hline
\texttt{-{}-{}} & en dash – & \texttt{\&ndash;}\\
\texttt{-{}-{}-} & em dash — & \texttt{\&mdash;}\\
\texttt{...} & ellipsis … & \texttt{\&hellip;}\\
\texttt{\textbackslash{}-} & soft hyphen & \texttt{\&shy;}\\
\end{tabular}
\end{center}

Smart quotes (\texttt{\#+OPTIONS: ':t}, off by default) turn "straight" quotes
into curly ones, using the quotes of \texttt{\#+LANGUAGE:}.

\textbf{Try:} export "Dashes and dots", then "Curly quotes".

\textbf{Expect:}
\begin{verbatim}
pages 10&ndash;20 &mdash; and so on&hellip;
\end{verbatim}
and for "Curly quotes":
\begin{verbatim}
&ldquo;Hello,&rdquo; she said. It&rsquo;s fine.
\end{verbatim}
\subsection{Dashes and dots}
\label{sec:orgcda24b4}
pages 10--20 --- and so on\ldots{}
\subsection{Curly quotes}
\label{sec:org14c171c}
"Hello," she said. It's fine.
\section{\LaTeX{} fragments}
\label{sec:org42f874c}
Math is written in \LaTeX{}: \texttt{\$x\textasciicircum{}2\$} or \texttt{\textbackslash{}(x\textasciicircum{}2\textbackslash{})} inline, \texttt{\textbackslash{}[ ... \textbackslash{}]} or
\texttt{\textbackslash{}begin\{equation\}} \ldots{} \texttt{\textbackslash{}end\{equation\}} for displayed formulas. In the
buffer they are highlighted (\texttt{OrgLatex}); \texttt{<prefix>xl} (Emacs
\texttt{<C-c><C-x><C-l>}) previews them as images if you have \LaTeX{} and an image
backend. HTML export uses MathJax. Details and previews:
\href{21-images-latex.org}{21-images-latex.org}.

\texttt{\$} is only math when it is not next to a space inside: \texttt{\$5 and \$10} is
money, not math.

\textbf{Try:} export "Math examples".

\textbf{Expect:}
\begin{verbatim}
The area is \(\pi r^2\) and \(a+b\).
\[ E = mc^2 \]
It costs $5 and $10.
\end{verbatim}
\subsection{Math examples}
\label{sec:org765f93b}
The area is \(\pi r^2\) and \(a+b\).
\[ E = mc^2 \]
It costs \$5 and \$10.
\section{Paragraphs, line breaks and rules}
\label{sec:org2ee9fa0}
\begin{itemize}
\item A \textbf{paragraph} is a group of lines; blank lines separate paragraphs. Line
breaks inside a paragraph are ignored when exporting.
\item \texttt{\textbackslash{}\textbackslash{}} at the end of a line forces a line break.
\item \texttt{\#+OPTIONS: \textbackslash{}n:t} keeps every line break of the file.
\item A line of \textbf{five or more dashes} (\texttt{-{}-{}-{}-{}-}) is a horizontal rule.
\end{itemize}

\textbf{Try:} export "Breaks and rules".

\textbf{Expect:} the first paragraph has no \texttt{<br />} (a browser shows its two
lines as one), the second has one after "Forced", and the dashes became
\texttt{<hr />}:
\begin{verbatim}
<p>
These two lines
form one paragraph.
</p>

<p>
Forced <br />
break here.
</p>

<hr />
\end{verbatim}
\subsection{Breaks and rules}
\label{sec:org4f04a7e}
These two lines
form one paragraph.

Forced \\
break here.

\noindent\rule{\textwidth}{0.5pt}

A new paragraph after the rule.
\section{Blocks}
\label{sec:orgd42c180}
Blocks start with \texttt{\#+begin\_NAME} and end with \texttt{\#+end\_NAME} (upper or lower
case). Insert them with \texttt{<prefix>ib} (see
\href{01-outline.org}{01-outline.org}). \texttt{<Tab>} on the \texttt{\#+begin\_} line
folds the block; \texttt{\#+STARTUP: hideblocks} folds all of them when the file
opens.

\begin{center}
\begin{tabular}{ll}
Block & Meaning\\
\hline
\texttt{quote} & a quotation\\
\texttt{center} & centered text\\
\texttt{verse} & a poem: line breaks and indentation are kept\\
\texttt{example} & verbatim text in a monospace font\\
\texttt{src LANG} & source code in language LANG\\
\texttt{export BACKEND} & raw text for one exporter only (\texttt{html}, \texttt{latex}\ldots{})\\
\texttt{comment} & never exported\\
\end{tabular}
\end{center}

Markup works inside quote, center and verse blocks, but not inside example
and src blocks.

\textbf{Try:} on "Block examples", press \texttt{<Tab>} until everything shows, then
press \texttt{<Tab>} on each \texttt{\#+begin\_} line to fold and unfold it.

\textbf{Expect:} each block folds into its \texttt{\#+begin\_} line.

\textbf{Try:} export "Block examples".

\textbf{Expect:} in the HTML you find \texttt{<blockquote>} with \texttt{<b>ideas</b>} inside,
\texttt{<div class}"org-center">=, \texttt{<p class}"verse">= with \texttt{\&nbsp;} for the
indentation and \texttt{<br />} at every line end, \texttt{<pre class}"example">= with
\texttt{*not bold*} left as it is, \texttt{<pre class}"src src-lua">=, the raw
\texttt{<em>raw HTML</em>}, and \textbf{no} trace of "This is a comment block".
\subsection{Block examples}
\label{sec:org15303e2}
\begin{quote}
Everything is made of \textbf{ideas}. --- Anonymous
\end{quote}

\begin{center}
Centered title
\end{center}

\begin{verse}
Roses are red,\\
\hspace*{2\fontdimen2\font}violets are blue,\\
\hspace*{4\fontdimen2\font}verses keep\\
\hspace*{6\fontdimen2\font}their indentation too.\\
\end{verse}

\begin{verbatim}
Example text: *not bold*, shown exactly as typed.
\end{verbatim}

\begin{verbatim}
print("hello from Lua")
\end{verbatim}
\subsection{Line numbers in examples}
\label{sec:org97668aa}
Example and src blocks accept switches. \texttt{-n} numbers the lines, \texttt{+n}
continues the numbering of the previous block.

\textbf{Try:} export "Numbered example".

\textbf{Expect:} the lines start with \texttt{<span class}"linenr">1: </span>= and
\texttt{<span class}"linenr">2: </span>=.
\subsection{Numbered example}
\label{sec:org8457a29}
\begin{verbatim}
1  first line
2  second line
\end{verbatim}
\section{Fixed-width lines}
\label{sec:orgd54f3f3}
A line starting with a colon and a space (=: =) is shown and exported
verbatim, like a one-line example block. Handy for short program output.

\texttt{<C-c>:} (\texttt{toggle\_fixed\_width}) adds or removes the =: = prefix on the
current line, or on every line of a Visual selection.

\textbf{Try:} select the two "output" lines below with \texttt{V} and \texttt{j}, press \texttt{<C-c>:}.

\textbf{Expect:} both lines start with \texttt{: = and are highlighted as fixed-width.
Select the two lines again (=V} and \texttt{j}) and press \texttt{<C-c>:} to remove it.

\textbf{Try:} export "Fixed-width example".

\textbf{Expect:}
\begin{verbatim}
<pre class="example">
$ date
Mon Sep 28 10:00:00 2026
</pre>
\end{verbatim}
\subsection{Fixed-width playground}
\label{sec:orgd81a00d}
output line one
output line two
\subsection{Fixed-width example}
\label{sec:orgbc77e6e}
\begin{verbatim}
$ date
Mon Sep 28 10:00:00 2026
\end{verbatim}
\section{Comments}
\label{sec:org8f24d88}
Three ways to keep text out of the export:

\begin{enumerate}
\item A line starting with \texttt{\#} followed by a space (or \texttt{\#} alone) is a
comment line. \texttt{\#+} starts a keyword instead, and \texttt{\#word} is plain text.
\item A \texttt{\#+begin\_comment} / \texttt{\#+end\_comment} block (see Blocks).
\item A headline starting with \texttt{COMMENT} comments out the whole subtree: it is
left out of the export and the agenda. \texttt{<prefix>hC} (Emacs \texttt{<C-c>;})
toggles the keyword.
\end{enumerate}

\textbf{Try:} export "Comment examples".

\textbf{Expect:} only "Visible text." and "\#hashtag is not a comment." are in the
HTML. Nothing of "Secret child" appears.

\textbf{Try:} on "Secret child", press \texttt{<prefix>hC}, and export again.

\textbf{Expect:} the keyword is removed, and "Secret child" and "Its text." now
appear in the HTML.
\subsection{Comment examples}
\label{sec:org73cb48e}
Visible text.

\#hashtag is not a comment.
\section{Keywords}
\label{sec:org1fec826}
Lines like \texttt{\#+TITLE: ...} are keywords: settings for the file or the
export. The ones at the top of this file are:

\begin{description}
\item[{\texttt{\#+TITLE:} and \texttt{\#+AUTHOR:}}] the document title and author.
\item[{\texttt{\#+STARTUP:}}] how the file opens (see
\href{01-outline.org}{01-outline.org}).
\item[{\texttt{\#+CATEGORY:}}] the name used by the agenda.
\item[{\texttt{\#+MACRO:}}] text macros (see Macros below).
\end{description}

Others you will meet: \texttt{\#+OPTIONS:} (export options), \texttt{\#+TODO:},
\texttt{\#+TAGS:}, \texttt{\#+PROPERTY:}, \texttt{\#+COLUMNS:}, and the \textbf{affiliated} keywords
\texttt{\#+NAME:}, \texttt{\#+CAPTION:} and \texttt{\#+ATTR\_HTML:} which go right above an element
(a table, a block, an image) to name or describe it. After editing a \texttt{\#+}
line, press \texttt{<C-c><C-c>} on it so the buffer picks up the change.
\section{Export snippets}
\label{sec:orgde2857d}
\texttt{@@BACKEND:text@@} inserts raw text for one exporter only, inline. Other
exporters drop it. Use it for things Org markup can't express, e.g.
keyboard keys in HTML.

\textbf{Try:} export "Snippet example" to HTML, then to UTF-8
(\texttt{:Org export utf8 subtree body buffer}).

\textbf{Expect:} in HTML: \texttt{Press <kbd>Ctrl</kbd>-<kbd>S</kbd> to save.} In UTF-8:
\texttt{Press Ctrl-S to save.}
\subsection{Snippet example}
\label{sec:org8242909}
Press Ctrl-S to save.
\section{Macros}
\label{sec:orgb72b707}
Macros are templates expanded when exporting: \texttt{\#+MACRO: name text} defines
one, \texttt{\{\{\{name(arg1,arg2)\}\}\}} uses it. In the text, \texttt{\$1}, \texttt{\$2}\ldots{} are the
arguments and \texttt{\$0} all of them. A comma inside an argument is written
\texttt{\textbackslash{},}. This file defines three macros at the top:

\begin{verbatim}
#+MACRO: greet Hello, $1!
#+MACRO: swap $2 before $1
#+MACRO: kbd @@html:<kbd>@@$1@@html:</kbd>@@
\end{verbatim}

Built-in macros:

\begin{center}
\begin{tabular}{ll}
Macro & Expands to\\
\hline
\texttt{\{\{\{title\}\}\}} & the \texttt{\#+TITLE:}\\
\texttt{\{\{\{author\}\}\}} & the \texttt{\#+AUTHOR:}\\
\texttt{\{\{\{date(FMT)\}\}\}} & the \texttt{\#+DATE:}, formatted (FMT optional)\\
\texttt{\{\{\{time(FMT)\}\}\}} & the export time, e.g. \texttt{\%Y-\%m-\%d}\\
\texttt{\{\{\{keyword(NAME)\}\}\}} & the value of any \texttt{\#+NAME:} keyword\\
\texttt{\{\{\{property(NAME)\}\}\}} & a property of the entry around the macro\\
\texttt{\{\{\{n\}\}\}} & a counter: 1, 2, 3\ldots{} (\texttt{\{\{\{n(name)\}\}\}})\\
\texttt{\{\{\{input-file\}\}\}} & the name of this file\\
\texttt{\{\{\{modification-time(FMT)\}\}\}} & when the file was last changed\\
\end{tabular}
\end{center}

In the buffer, macros are highlighted with \texttt{OrgMacro} and not expanded.

\textbf{Try:} export "Macro examples".

\textbf{Expect:} (the export also has a small table of contents, because of the
child headline)
\begin{verbatim}
Hello, World! Hello, Ada, Grace!
second before first
Press <kbd>C-c</kbd>.
Title: Markup: emphasis, blocks, comments and friends
Author: Ada Lovelace
Steps 1, 2, 3.
\end{verbatim}
and under the "Weekly meeting" heading:
\begin{verbatim}
Room: Blue room
\end{verbatim}
\subsection{Macro examples}
\label{sec:orgd57504c}
Hello, World! Hello, Ada, Grace!
second before first
Press C-c.
Title: Markup: emphasis, blocks, comments and friends
Author: Ada Lovelace
Steps 1, 2, 3.
\subsubsection{Weekly meeting}
\label{sec:org0c92d11}
Room: Blue room
\section{Further reading}
\label{sec:orgb0f1a4e}
\begin{description}
\item[{\texttt{:h org-appearance}}] hiding markers, pretty entities, highlights
\item[{\texttt{:h org-highlights}}] the highlight groups (\texttt{OrgBold}, \texttt{OrgLatex}\ldots{})
\item[{\texttt{:h org-structure}}] \texttt{<prefix>E}, \texttt{<prefix>ib}, \texttt{<C-c>:}
\item[{\texttt{:h org-export-settings}}] \texttt{\#+OPTIONS:} keys and macros
\item[{\texttt{:h org-images}}] \LaTeX{} previews
\item \href{19-export.org}{19-export.org} and
\href{21-images-latex.org}{21-images-latex.org} :: more on export and math
\end{description}
